% Requisitos ampliados

\subsection{Ejercicio 7}\label{subsec:ej7}
La función \texttt{bwfilter()} aplica un filtro de frecuencia, eliminando las 
muestras entre las frecuencias (en Hz) indicadas.
\begin{lstlisting}[caption={Filtro de frecuencia}, label={lst:filtro}]
filtro <- bwfilter(basico, f=basico@samp.rate, from=10000, to=20000, bandpass= FALSE, output = "Wave")

writeWave(filtro, file.path("filtrado.wav"))
\end{lstlisting}

Al poner \texttt{f = basico@samp.rate} la función lee automáticamente la cabecera del archivo y aplica el filtro de frecuencias de manera exacta y sin errores.

\begin{figure}[ht]
    \centering
    \includegraphics[width=0.95\linewidth]{images/7_filtro.png}
    \caption{Filtro de frecuencia}
    \label{fig:filtro}
\end{figure}

En la Figura \ref{fig:filtro} da un aviso al ejecutar la función \texttt{writeWave}. Este aviso no afecta al ejercicio ya que el redondeo no es perceptible para el oído humano.\\

Para ver información sobre el sonido uso la función \texttt{str} y la propia llamada al sonido. Se puede observar en \textit{Listing} \ref{lst:filtroinfo} y \textit{Figura} \ref{fig:filtroinfo}.
\begin{lstlisting}[caption={Información del filtro de frecuencia}, label={lst:filtroinfo}]
str(filtro)
filtro
\end{lstlisting}

\begin{figure}[ht]
    \centering
    \includegraphics[width=0.95\linewidth]{images/7_2infoFiltro.png}
    \caption{Información del \textit{filtro}}
    \label{fig:filtroinfo}
\end{figure}
%----------------------------------------------------------------

\subsection{Ejercicio 8}\label{subsec:ej8}
